\documentclass[11pt]{article}
\usepackage[T1]{fontenc}
\usepackage{amsmath,amssymb,hyperref,geometry}
\geometry{margin=1in}
\title{Native Feedback and Interaction Extension\\Research Contract v0.2}
\author{}\date{16 November 2026}
\begin{document}\maketitle

\paragraph{Purpose and status.}
Investigate whether the original category-valued relative context/content construction admits a conservative native extension that systematically generates and relates admissible feedback, dependent occurrence substitution and shared-effect interaction. Discovery and characterization of that larger structure are the governing objective, not merely the addition of one coherent operation. A single completion, or a coherent family of extensions with proved relationships, may supply the answer. Existence of either is a research question, not an assumption.

This revision replaces v0.1's branching-first default and selected-investigation completion criterion while retaining its structured-input, native-witness and conservativity safeguards. It does not invalidate existing results, redefine completion of the feedback-and-concurrency contract v0.1.3, or require a contract-closure audit. The earlier contract's parameterized completion criterion is not a ceiling on the mathematical generality sought here.

\section{Governing question and research scope}
\begin{quote}
\textbf{Does the original relative context/content construction admit a conservative native completion that provides a structural account of feedback and shared interaction, including their relationship to dependent and branching graded occurrences? What construction, universal property or precise obstruction determines its domain and limits?}
\end{quote}

A free, reflective, relative, coinductive or other universal characterization is a preferred possibility, not a prescribed architecture or promised theorem. Investigate whether existing state-feedback, native-grade, interaction and guarded-prefix constructions arise as restrictions, instances or stages of such a structure. A coherent family of constructions is acceptable when its assembly and compatibility are proved; a catalogue of unrelated theorems is not a structural answer.

General admissible feedback remains a central research obligation. Native branching with interference and infinite semantic completion are important directions and discriminating tests, not substitutes for that obligation. None is prescribed as the first route. Their actual dependencies must be discovered. They need not be solved simultaneously, and no unrestricted total trace, largest completion over unspecified doctrines, or universal infinite-grade object is demanded.

Generality must be indexed by declared source doctrines, feedback connections, occurrence shapes, interaction wirings and preservation judgments. Seek reusable structural hypotheses explaining the admitted classes and their boundaries, rather than making the class coincide by definition with examples already known to work. An unresolved direction remains open; it cannot be removed from the governing question merely because an easier extension succeeds.

\section{Baseline and structured inputs}
Begin with the independently specified context and content presentations, pure interface, native grade maps, changing-endpoint body squares and coherent crossing underlying
\[
 M(A,B)=(e_A\downarrow B),\qquad
 N(A,B)=(A\downarrow d_B),\qquad
 H(A,B)=(e_A\downarrow d_B).
\]
The underlying hom categories alone are not the input. State which admission diagrams, joint effects, coupling choices, reindexings, comparison cells and coherence data each proposed extension consumes. Supplied choices remain explicit data unless a proved characterization determines them up to the stated equivalence. Distinct native maps must not be identified merely because a carrier interpretation cannot distinguish them.

Compare at least the original finitary polynomial-occurrence construction and the newer local-joint construction for independently graded one-hole content. The former already supports branching under its hypotheses; the latter supplies shared-effect comparisons and native coefficient lifts. Earlier addressed branching uses an independence policy that is not interchangeable with an interfering coupling. Thus neither ``branching exists'' nor ``a recursive runtime exists'' establishes the proposed extension. Native multi-occurrence substitution with chosen shared effects is a significant unresolved combination at the v0.1 baseline, not a prior determination of the only missing structure.

Observable-cover and residual-profile certificates address declared tests, readiness and information. They do not replace actual native witnesses, probability laws or grade arrows. Existing finite-support, infinite-grade, state-copying and coherent-lifting obstructions constrain candidate constructions; repeating one does not constitute a new result. Pin the source versions used and recheck these baseline boundaries when later results change them.

\section{Research obligations}

\subsection*{E0: Comparative structural diagnosis}
Identify the missing capabilities at their actual interfaces. Distinguish nonexpressibility, missing native cells, joint admission, failed coherence, absent semantic interpretation and absence of a structural characterization. Do not infer impossibility from an unproved theorem or an unavailable implementation.

Compare leading completion mechanisms against the same mathematical probes before selecting a route. Feedback/iteration-first, occurrence/interaction-first and constructions generating both are candidates, not mandatory architectures. Use existing constructions and primary literature where they discharge the obligations. Record the independent inputs, expected outputs and first unresolved compatibility problem for each live route. Justify pruning by a theorem, counterexample or concrete mismatch with a required interface, not convenience alone. Derive the dependency structure rather than imposing a branching-first sequence.

\subsection*{E1: Structured extension and characterization}
Specify the structured source and target doctrines and their structure-preserving maps, including the comparison cells required at the claimed categorical dimension. Construct the proposed extension and its inclusion or interpretation of the old presentation. State the functoriality or higher coherence actually needed. Distinguish formally adjoining an operation from interpreting it in a substantive semantic class.

Investigate an appropriate mapping, extension, freeness, reflection or other universal property. Define its target class and maps independently of the desired conclusion, and prove the relevant existence, uniqueness or equivalence and its coherence. A syntactic universal property does not by itself supply semantic existence or conservation of native distinctions. Conversely, failure of one proposed universal property does not rule out every useful extension.

If a single universal completion is unavailable or inappropriate, give a structural alternative: for example, compatible relative constructions with proved comparison and assembly results. Explain why this alternative answers the governing question for the declared class. The universal-characterization question must receive a proved answer, a precisely scoped obstruction, or an explicit unresolved status; it must not disappear after a local construction succeeds.

\subsection*{E2: Feedback, iteration and admissible interconnection}
Distinguish retained-state closure, seeded delayed feedback, guarded recursive control, instantaneous cyclic interconnection and general iteration. Determine which can be generated, interpreted or related by the proposed extension, and which require genuinely independent structure. The current guarded regime is a baseline, not an a priori upper bound on the investigation.

For each claimed feedback operation specify its typed domain of admissible connections, action on native cells, and behavior under reindexing, dependent substitution and joint interaction. Prove the applicable unfolding, unit, naturality, nesting, sliding or iteration laws at their actual equality or directed-comparison strength. Do not import a trace axiom inappropriate to the selected semantics or invert an effectful comparison without proof. Fixed-point existence, uniqueness, selection and preservation by representation changes are separate obligations.

Derive admissibility and comparison preservation from independent structural hypotheses. A whole closed-system comparison square, or the assertion that feedback already preserves every desired law, is not an acceptable replacement for that derivation. Partial, relational, guarded and other restricted constructions may be legitimate, but their domains and failures must be explained structurally. Determine what extends beyond the old feedback hypotheses, or what obstruction prevents the proposed enlargement. Merely retaining the old feedback operation beside a new branching operation leaves this question open unless their structural relationship is established.

\subsection*{E3: Dependent occurrences and shared interference}
Specify occurrence shapes, changing grades and interfaces, result-dependent continuations, joint admission and retained effects. Investigate native multi-occurrence substitution with chosen shared effects, including its unit, associativity, reindexing, comparison and interaction coherence. Primitive grade and boundary data must construct actual native cells; pointwise existence of fillers is not coherent selection.

Keep the interaction boundary explicit. Do not duplicate a shared state into independent copies, infer a joint effect from component marginals, replace a chosen coupling by a product law, or assume serializability. Retain intermediate post-contexts, dependencies and pending effects whenever the intended judgments consume them. State how admission and coupling choices transform under substitution and regrouping, rather than independently choosing incompatible joints at each stage.

Relate the resulting occurrence and interaction operations to the feedback structures of E2 through proved compatibility maps and equations. Neither the existence of each operation separately nor a unary encoding alone supplies this relationship. An encoding may contribute only with the required interpretation, native distinctions and structural laws proved; an already available encoding without new categorical capability is not the sought extension.

\subsection*{E4: Recurrent behavior and semantic completion}
Separate infinite raw histories, coherent finite native comparisons, compatible finite effect laws, divergence, probability laws on infinite executions and infinite coefficient or content objects. Identify which are actually consumed by each proposed feedback or interaction operation. Separate witnesses at every horizon do not establish one coherent unbounded comparison.

When infinite objects are needed, construct an appropriate completion or prove an obstruction at the specified interface. Verify the required existence, transport, continuity or limit-preservation hypotheses in the actual native presentation. Explain any change of ambient category and its relation to E5. Do not assume that completing effects also completes grades, or that adding either completion preserves the mixed crossing and joint admission automatically.

A coherent prefix or coinductive semantics may be sufficient without an accumulated infinite grade. Prove that adequacy for the claimed operation instead of introducing unnecessary limits. Conversely, lack of an immediate need for limits in a finite example does not settle the broader completion question. Record whether each completion issue is resolved, structurally unnecessary for a stated operation, obstructed, or open.

\subsection*{E5: Conservativity and recovery}
Declare exactly which old objects, native arrows, equations, observations and operations the extension preserves and reflects. Specify whether new comparisons between old objects are permitted. Prove the selected faithfulness, equation-reflection or other conservativity statement; local fullness is not implicit in the word ``conservative.'' Preserve the independently presented context/content structure and crossing through explicit comparison, rather than retaining only their carrier-level composite.

Recover the appropriate original finite, independent and guarded cases with their actual hypotheses. Explain when a new observer separates an old quotient without invalidating its old theorem at its old observation. Added operations must not silently equate different native grade maps, discard retained coupling choices or identify refusal, deadlock and divergence. A necessary change of equality, observation or representation must be declared and justified, not described as unchanged recovery.

\subsection*{E6: Strict extension and substantive interpretation}
Exhibit a genuinely new process, operation, admissible interconnection or compositional theorem not supplied by the pinned baseline under its established hypotheses. Classify the separation as nondefinability, failure of interpretation, or nonderivability from those hypotheses, and provide evidence at that strength. No claim of nondefinability follows merely from the absence of an existing construction.

Supply a substantive interpretation of the new structure, preferably combining nontrivial native grades, branching recurrence and shared interference in the same development. Test competing constructions on common probes. Disconnected examples may clarify independent obligations but do not alone establish their composition. The semantic interpretation must verify the primitive admission, effect, grade and body hypotheses; a free syntax or change of notation alone is insufficient.

Use obstructions to delimit or improve the construction. State the exact class ruled out, identify the violated obligation, and construct a meaningful repaired alternative where possible. A restricted alternative is not automatically strongest or optimal. Old counterexamples and a failed candidate are not substitutes for a new structural result, nor do they imply that all completions fail.

\section{Research sequence and completion evidence}
The obligations describe proof dependencies and acceptance conditions, not a fixed execution order. A first branching theorem, a new feedback fragment or a semantic completion can each be useful intermediate progress. Record how it advances or discriminates the governing structural question. Do not make arbitrary instantaneous feedback or infinite completion prerequisites for every intermediate result; do not turn their postponement into their removal from the investigation.

\paragraph{Positive structural outcome.}
A positive answer at a declared meaningful scope supplies the structured extension problem; a construction or compatible family with proved native coherence; the feedback/interaction relationship; the conservativity and recovery maps; a structural characterization; and a separating semantic interpretation. Explain how the existing constructions fit within, map into, or fail to extend to the new structure. Characterization may be universal or a justified structural alternative under E1, but cannot consist only of placing operations side by side.

The conclusion must account for general feedback, graded occurrence interaction and recurrent/infinite semantics. An account may establish a construction, derive its relationship to another, prove a relevant obstruction, or justify why a particular completion is unnecessary. It need not supply all operations in every doctrine. If a load-bearing part of the stated governing question remains unaddressed, report a partial answer rather than redefine the question to match the result. In particular, one new coherent branching operation compatible only with already supplied feedback is an intermediate result unless the broader structural characterization is also established.

\paragraph{Obstruction or restricted structural outcome.}
A precise no-go theorem can resolve a declared completion problem, and a repaired construction may supply a conditional structural answer. Relate the obstruction to the competing routes and to the governing question. It completes only the class of problems it actually rules out; failure of one route is not closure of the general investigation. No classification of all possible doctrines is required. Equally, relabelling a convenient fragment as the whole target is not a completed structural resolution.

\paragraph{Progress and unresolved outcomes.}
Keep distinct the status of intermediate theorems, the characterization of a candidate extension, and the answer to the governing question. A count of new manuscripts, successful finite tests or additional applications cannot substitute for the missing assembly or characterization theorem. If neither a positive nor a relevant obstruction outcome is obtained, identify the remaining mathematical obligation without simulating completion. Scope revisions require an explicit mathematical reason and retain the displaced question as open.

\section{Evidence and source discipline}
Separate established imports, adaptations and potentially new claims. Verify primary-source hypotheses and the interpretation that makes each import applicable. A candidate construction earns additional higher structure through a named operation, universal property or proof obligation, not merely because that structure is a natural next abstraction. Obtain independent review of load-bearing written arguments and record its scope; distinguish written proof, finite diagnostics, formal proof and source checks.

This is not a contract-closure audit, an application-coverage exercise, or a requirement to finish an external Byzantine-control mission. New applications may test a categorical claim but cannot replace it. No Lean installation, formalization campaign or TeX compilation is required. Preserve source checkpoints and reviews outside the transient workspace; do not claim independent review or mechanized verification before it occurs.

\paragraph{Authoritative starting sources.}
Paths are relative to the contextual-composition project:
\begin{quote}\small\raggedright
\path{synthesis/Integrated_Relative_Computation_Acceptance_Dossier.tex}\par
\path{foundations/local-joint-comparison/Local_Joint_Comparison_and_Generated_Interference.tex}\par
\path{feedback-and-concurrency/observable-covers/Observable_Covers_and_Local_Observers.tex}\par
\path{feedback-and-concurrency/observable-covers/Residual_Profiles_and_Admission_Reflection.tex}\par
\path{contracts/Feedback_and_Concurrency_Research_Contract_v0_1_3.tex}\par
\path{contracts/Native_Feedback_and_Interaction_Extension_Contract_v0_1.tex}
\end{quote}
Pin the versions used by each mathematical result. This contract repairs the research target; it does not supply permanent assumptions about later source changes or certify the current research results. Further authoritative dependencies should be added when a proposed construction actually consumes them.

\end{document}
